
\subsection{Temperature and density}

In order to use parameters without dimensions in the calculations, all of
them are normalized to reference values. Here temperature and densities
correspond to flux surface-averaged values, though the usual bracket
notation is not introduce, for sake on simplicity.  

Normalized temperatures, are defined by the relations 
\begin{equation}
\begin{array}{l}
\overline{T}_{e}\left( \psi \right) =T_{e}\left( \psi \right) /T_{e}^{\dag }
\\ 
\overline{T}_{ss^{\prime }}\left( \psi \right) =T_{ss^{\prime }}\left( \psi
\right) /T_{e}^{\dag }
\end{array}
\end{equation}
where the subscripts $e$ stands for electrons and $ss^{\prime }$ for ion
species $s$ in the ionization state $s^{\prime }$, whose charge is $%
Z_{ss^{\prime }}$. Here, both $T_{e}\left( \psi \right) $ and $T_{e}^{\dag }$
have defined units, while $T_{e}^{\dag }$ is the reference value for the
normalization. According to this rule, the normalized densities are defined
in the same way 
\begin{equation}
\left\{ 
\begin{array}{l}
\overline{n}_{e}\left( \psi \right) =n_{e}\left( \psi \right) /n_{e}^{\dag }
\\ 
\overline{n}_{ss^{\prime }}\left( \psi \right) =n_{ss^{\prime }}\left( \psi
\right) /n_{e}^{\dag }
\end{array}
\right. 
\end{equation}
and in order to ensure that all quantities are less than unity from the
plasma center to the edge, references values are given by the condition 
\begin{equation}
\left\{ 
\begin{array}{l}
T_{e}^{\dag }=\max \left( T_{e}\left( \psi \right) ,T_{ss^{\prime }}\left(
\psi \right) \right)  \\ 
n_{e}^{\dag }=\max \left( n_{e}\left( \psi \right) ,n_{ss^{\prime }}\left(
\psi \right) \right) 
\end{array}
\right. 
\end{equation}
for $0\leq \psi \leq \psi _{\max }.$

In principle, ion densities $n_{ss^{\prime }}\left( \psi \right) $ are
obtained from a particle transport code transport code. However, when the
plasma is made of two types of ions which are fully ionized (low $Z$ limit),
a simple model may be used in order to determine densities $n_{s}\left( \psi
\right) $ with $s=\left\{ 1,2\right\} $, knowing electron density $%
n_{e}\left( \psi \right) $ and effective charge $Z_{eff}\left( \psi \right) $
profiles. Indeed, since by definition, 
\begin{eqnarray}
Z_{eff}\left( \psi \right)  &=&\sum_{ss^{\prime }}\frac{n_{ss^{\prime
}}\left( \psi \right) Z_{ss^{\prime }}^{2}}{n_{ss^{\prime }}\left( \psi
\right) Z_{ss^{\prime }}}  \nonumber \\
&=&\frac{1}{n_{e}\left( \psi \right) }\sum_{ss^{\prime }}n_{ss^{\prime
}}\left( \psi \right) Z_{ss^{\prime }}^{2}
\end{eqnarray}
taking into account of the local electro-neutrality, $n_{e}\left( \psi
\right) =\sum_{ss^{\prime }}n_{ss^{\prime }}\left( \psi \right)
Z_{ss^{\prime }},$ one obtains in this simple case a system of two equations
with two unknowns, 
\begin{equation}
\left\{ 
\begin{array}{l}
n_{1}\left( \psi \right) Z_{1}^{2}+n_{2}\left( \psi \right)
Z_{2}^{2}=n_{e}\left( \psi \right) Z_{eff}\left( \psi \right)  \\ 
n_{1}\left( \psi \right) Z_{1}+n_{2}\left( \psi \right) Z_{2}=n_{e}\left(
\psi \right) 
\end{array}
\right. 
\end{equation}

Therefore, 
\begin{equation}
n_{1}\left( \psi \right) =n_{e}\left( \psi \right) \frac{\left(
Z_{eff}\left( \psi \right) -Z_{2}\right) }{Z_{1}\left( Z_{1}-Z_{2}\right) }
\end{equation}
\begin{equation}
n_{2}\left( \psi \right) =n_{e}\left( \psi \right) \frac{\left(
Z_{1}-Z_{eff}\left( \psi \right) \right) }{Z_{2}\left( Z_{1}-Z_{2}\right) }
\end{equation}
and using the normalization rule, 
\begin{equation}
\overline{n}_{1}\left( \psi \right) =\overline{n}_{e}\left( \psi \right) 
\frac{\left( Z_{eff}\left( \psi \right) -Z_{2}\right) }{Z_{1}\left(
Z_{1}-Z_{2}\right) }
\end{equation}
\begin{equation}
\overline{n}_{2}\left( \psi \right) =\overline{n}_{e}\left( \psi \right) 
\frac{\left( Z_{1}-Z_{eff}\left( \psi \right) \right) }{Z_{2}\left(
Z_{1}-Z_{2}\right) }
\end{equation}
since by definition $Z_{eff}\left( \psi \right) $ has no units.

In the calculations, one may also consider several isotopes, since ion
inertia play a significant role for the collision operator. Hence, 
\begin{equation}
n_{ss^{\prime }}\left( \psi \right) =\sum_{m}n_{ss^{\prime }}^{m}\left( \psi
\right) 
\end{equation}
or 
\begin{equation}
\overline{n}_{ss^{\prime }}\left( \psi \right) =\sum_{m}\overline{n}%
_{ss^{\prime }}^{m}\left( \psi \right) 
\end{equation}
where the masses of isotope $m$ of ion species $ss^{\prime }$ are $%
m_{ss^{\prime }}^{m}.$ In the calculations 
\begin{equation}
\overline{m}_{ss^{\prime }}^{k}=m_{ss^{\prime }}^{m}/m_{e}
\end{equation}
where $m_{e}$ is the electron rest mass.

\subsection{Time}

Regarding time ordering which is considered in the derivation of the drift
kinetic equation and the dominant role played by collisions for the electron
momentum distribution function build-up, it is natural to scale time
evolution to the reference electron thermal collision frequency $\nu
_{e}^{\dag }$ according to the relation 
\begin{equation}
\overline{t}=t/\tau _{c}^{\dag }=t\nu _{e}^{\dag }
\end{equation}
where 
\begin{equation}
\nu _{e}^{\dag }=\frac{q_{e}^{4}n_{e}^{\dag }\ln \Lambda ^{\dag }}{4\pi %
\varepsilon _{0}m_{e}^{2}c^{3}\beta _{th}^{\dag 3}}
\end{equation}

Here $\varepsilon _{0}=8.854187818\times 10^{-12}$ $\left( F/m\right) $ is
the free space permeability, $c=299792458$ $\left( km/s\right) $ the speed
of light, $q_{e}=-1.60217733\times 10^{-19}$ $\left( C\right) $ the electric
charge of the electron, $m_{e}=9.1093897\times 10^{-31}$ $\left( Kg\right) $
the electron rest mass, and 
\begin{equation}
\ln \Lambda ^{\dag }=25.2-0.5*\log (n_{e}^{*})+\log (T_{e}^{*})
\end{equation}
is the usual reference Coulomb logarithm, with 
\begin{equation}
\beta _{th}^{\dag }=\sqrt{\frac{T_{e}^{\dag }}{m_{e}c^{2}}}=\sqrt{\Theta
^{\dag }}
\end{equation}
where $m_{e}c^{2}=510.99905$ $\left( keV\right) $is the electron rest mass
energy.

\subsection{Momentum, velocity, and kinetic energy}

In the calculations, the momentum $p$ in relativistic units $\left(
m_{e}c\right) $ is normalized to the thermal reference value $p_{th}^{\dag
}=m_{e}v_{th}^{\dag },$%
\begin{equation}
\overline{p}=p/p_{th}^{\dag }
\end{equation}
and consequently
\begin{equation}
p_{th}^{\dag }/m_{e}c\approx v_{th}^{\dag }/c=\beta _{th}^{\dag }
\end{equation}
since thermal electrons are only weakly relativistic. The well known
relativistic correction factor $\gamma $ is then simply given by the relation
\begin{equation}
\gamma =\sqrt{p^{2}+1}=\sqrt{\overline{p}^{2}\beta _{th}^{\dag 2}+1}
\end{equation}
and in the non-relativistic limit, i.e. when $\overline{p}^{2}\beta
_{th}^{\dag 2}\ll 1,$ $\gamma \approx 1.$

Since in relativistic units, the total energy is linked to the relativistic
momentum according to the expression,
\begin{equation}
\left( E_{c}+1\right) ^{2}=\gamma 
\end{equation}
it is straightforward to expression the kinetic energy $E_{c}$ as a function
of  $\overline{p}$ in units of electron rest mass energy $m_{e}c^{2}$%
\begin{equation}
E_{c}=m_{e}c^{2}\left( \sqrt{\overline{p}^{2}\beta _{th}^{\dag 2}+1}%
-1\right) 
\end{equation}

Finally, concerning the normalization of the electron velocity $v$, one has
\begin{eqnarray}
v/c &=&p/\left( \gamma m_{e}c\right)   \nonumber \\
&=&\overline{p}p_{th}^{\dag }/\left( \gamma m_{e}c\right)   \nonumber \\
&=&\overline{p}\beta _{th}^{\dag }/\gamma 
\end{eqnarray}
and using $\overline{v}=v/v_{th}^{\dag },$ it comes
\begin{equation}
\overline{v}v_{th}^{\dag }/c=\overline{p}\beta _{th}^{\dag }/\gamma 
\end{equation}
or
\begin{equation}
\overline{v}=\overline{p}/\gamma 
\end{equation}
with 
\begin{equation}
v_{th}^{\dag }/c=\beta _{th}^{\dag }
\end{equation}

\subsection{Maxwellian electron momentum distribution}

With the definition of $\overline{p}$, the approximate relativistic
Maxwellian electron momentum distribution function $f_{M}\left( \psi
,p\right) $ up to first order $\gamma $ correction used in the code is given
by
\begin{equation}
f_{M}\left( \psi ,p\right) \simeq \frac{n_{e}\left( \psi \right) }{\left[ 2%
\pi T_{e}\left( \psi \right) \right] ^{3/2}}\exp \left[ -\frac{p^{2}}{\left(
1+\gamma \right) T_{e}\left( \psi \right) }\right] 
\end{equation}
where $T_{e}\left( \psi \right) $ is expressed in electron rest mass energy $%
m_{e}c^{2}$ unit, and providing that $\gamma -1\ll 1.$ Therefore, the
condition
\begin{equation}
\overline{p}^{2}\beta _{th}^{\dag 2}\ll 1
\end{equation}
must be fullfiled. Since $\overline{p}$ may be as large as $30$ in
calculations, in order to correctly describe momentum dynamics of the
fastest electrons, it results that 
\begin{equation}
\beta _{th}^{\dag }=\Theta ^{\dag 2}\ll 1/30
\end{equation}
or
\begin{equation}
T_{e}^{\dag }\ll 511\sqrt{1/30}\approx 100keV
\end{equation}

Usually, this condition is well fullfiled even in the core of tokamak
plasmas, where $T_{e}^{\dag }$ never exceeds $20keV.$ In normalized units, 

\begin{equation}
\overline{f}_{M}\left( \psi ,p\right) =\frac{\overline{n}_{e}\left( \psi
\right) n_{e}^{\dag }}{\left[ 2\pi T_{e}\left( \psi \right) \right] ^{3/2}}%
\exp \left[ -\frac{\overline{p}^{2}\beta _{th}^{\dag 2}}{\left( 1+\gamma
\right) T_{e}\left( \psi \right) }\right] 
\end{equation}
and since $\beta _{th}^{\dag }=\sqrt{T_{e}^{\dag }/m_{e}c^{2}},$ it turns
out that
\begin{equation}
\overline{f}_{M}\left( \psi ,p\right) \approx \frac{\overline{n}_{e}\left(
\psi \right) }{\left[ 2\pi \overline{T}_{e}\left( \psi \right) \right] ^{3/2}%
}\exp \left[ -\frac{\overline{p}^{2}}{\left( 1+\gamma \right) \overline{T}%
_{e}\left( \psi \right) }\right] 
\end{equation}
since
\begin{eqnarray*}
2\pi \int_{-1}^{+1}d\xi _{0}\int_{0}^{+\infty }\overline{p}^{2}\overline{f}%
_{M}\left( \psi ,p\right) d\overline{p} &=& \\
&&4\pi \frac{\overline{n}_{e}\left( \psi \right) }{\left[ 2\pi \overline{T}%
_{e}\left( \psi \right) \right] ^{3/2}}\int_{0}^{+\infty }\overline{p}%
^{2}\exp \left[ -\frac{\overline{p}^{2}}{\left( 1+\gamma \right) \overline{T}%
_{e}\left( \psi \right) }\right] d\overline{p}
\end{eqnarray*}

\subsection{Poloidal flux coordinate}

For the poloidal flux coordinate, the normalized value is 
\begin{equation}
\overline{\psi }=\psi /\psi _{a}
\end{equation}
where $\psi _{a}$ is the value at the plasma edge, corresponding usually to
the last closed magnetic flux surface, as given by the equilibrium code.
Usually, it is taken so that $0\leq \overline{\psi }\leq 1$ from the center
to the plasma edge.

\subsection{Drift kinetic coefficient}

The first order drift kinetic equation requires to calculate $\widetilde{f}$
defined as 
\begin{equation}
\widetilde{f}=\frac{v_{\Vert }I\left( \psi \right) }{\Omega }\frac{\partial %
f_{0}}{\partial \psi }
\end{equation}
where by definition $B_{T}\left( \psi \right) =\left| I\left( \psi \right)
\right| /R$ and $\Omega =-\frac{\left| e\right| B}{\gamma m_{e}}$. Since 
\begin{equation}
\overline{v}v_{th}^{\dag }/c=\overline{p}\beta _{th}^{\dag }/\gamma 
\end{equation}
one obtains in normalized units
\begin{eqnarray}
\overline{\widetilde{f}} &=&v\xi \frac{\gamma m_{e}}{q_{e}}R\frac{%
B_{T}\left( \psi \right) }{B}\frac{\partial \overline{f}_{0}}{\partial \psi }
\nonumber \\
&=&\overline{v}v_{th}^{\dag }\xi \frac{\gamma m_{e}}{q_{e}}\frac{B_{T}\left(
\psi \right) }{B}\frac{R}{\psi _{a}}\frac{\partial \overline{f}_{0}}{%
\partial \overline{\psi }}  \nonumber \\
&=&\overline{p}\xi \frac{cm_{e}}{q_{e}}\beta _{th}^{\dag }\frac{B_{T}\left(
\psi \right) }{B}\frac{R}{\psi _{a}}\frac{\partial \overline{f}_{0}}{%
\partial \overline{\psi }}
\end{eqnarray}

Therefore,
\begin{equation}
\overline{\widetilde{f}}=\overline{p}\xi \widetilde{C}\frac{\partial 
\overline{f}_{0}}{\partial \overline{\psi }}
\end{equation}
where
\begin{equation}
\widetilde{C}=\frac{cm_{e}}{q_{e}}\beta _{th}^{\dag }\frac{B_{T}\left( \psi
\right) }{B}\frac{R}{\psi _{a}}
\end{equation}

It is important to notice that values for calculating $\widetilde{C}$ must
be determined on the location where $B$ is minimum, according to the bounce
average procedure. 

For circular concentric flux surfaces, since $\partial /\partial \psi
=\left( 1/\nabla \psi \right) \partial /\partial r$ and $\,\nabla \psi
=RB_{P},$ it turns out that 
\begin{equation}
\widetilde{f}=v\xi \frac{\gamma m_{e}}{q_{e}}\frac{B_{T}\left( \psi \right) 
}{B_{P}}\frac{1}{B}\frac{\partial f_{0}}{\partial r}
\end{equation}
and naturally

\begin{equation}
\widetilde{C}\approx \frac{cm_{e}}{q_{e}}\beta _{th}^{\dag }\frac{q_{cyl.}}{%
a_{p}\epsilon }\frac{1}{B}
\end{equation}
where the safety factor $q$ is approximated by the cylindrical expression $%
q_{cyl.}=\frac{rB_{T}\left( \psi \right) }{RB_{P}}$, $a_{p}$ the plasma
minor radius, and $\epsilon =r/R_{0}$ the inverse aspect ratio.

\subsection{Momentum convection and diffusion}

From the conservative form of the Fokker-Planck equation in momentum space
\begin{equation}
\frac{\partial f_{0}}{\partial t}\propto \frac{\partial }{\partial p}\left(
S_{\mathbf{p}}\right) 
\end{equation}
where
\begin{equation}
S_{\mathbf{p}}\propto -\Bbb{D}_{\mathbf{p}}\frac{\partial f_{0}}{\partial p}+%
\mathbf{F}_{\mathbf{p}}f_{0}
\end{equation}

Introducing normalized coordinates, it turns out that
\begin{eqnarray}
\frac{\partial \overline{f}_{0}}{\partial \overline{t}} &\propto &\frac{1}{%
\nu _{e}^{\dag }p_{th}^{\dag }}\frac{\partial }{\partial \overline{p}}\left(
-\frac{\Bbb{D}_{\mathbf{p}}}{p_{th}^{\dag }}\frac{\partial \overline{f}_{0}}{%
\partial \overline{p}}+\mathbf{F}_{\mathbf{p}}\overline{f}_{0}\right)  
\nonumber \\
&=&\frac{\partial }{\partial \overline{p}}\left( -\frac{\Bbb{D}_{\mathbf{p}}%
}{\nu _{e}^{\dag }p_{th}^{\dag 2}}\frac{\partial \overline{f}_{0}}{\partial 
\overline{p}}+\frac{\mathbf{F}_{\mathbf{p}}}{\nu _{e}^{\dag }p_{th}^{\dag }}%
\overline{f}_{0}\right)   \nonumber \\
&=&\frac{\partial }{\partial \overline{p}}\left( -\overline{\Bbb{D}}_{%
\mathbf{p}}\frac{\partial \overline{f}_{0}}{\partial \overline{p}}+\overline{%
\mathbf{F}}_{\mathbf{p}}\overline{f}_{0}\right) 
\end{eqnarray}
where normalized diffusive and convective coefficients are
\begin{eqnarray}
\overline{\Bbb{D}}_{\mathbf{p}} &=&\Bbb{D}_{\mathbf{p}}/\Bbb{D}_{\mathbf{p}%
}^{\dag } \\
\overline{\mathbf{F}}_{\mathbf{p}} &=&\mathbf{F}_{\mathbf{p}}/\mathbf{F}_{%
\mathbf{p}}^{\dag }
\end{eqnarray}
with
\begin{eqnarray}
\Bbb{D}_{\mathbf{p}}^{\dag } &=&\nu _{e}^{\dag }p_{th}^{\dag 2} \\
\mathbf{F}_{\mathbf{p}}^{\dag } &=&\nu _{e}^{\dag }p_{th}^{\dag }
\end{eqnarray}

For the case of the Ohmic electric field, this normalization leads to
introduce naturally the Dreicer field
\begin{eqnarray}
E^{\dag } &=&\nu _{e}^{\dag }p_{th}^{\dag }/q_{e} \\
&=&m_{e}c\beta _{th}^{\dag }\nu _{e}^{\dag }/q_{e}
\end{eqnarray}
the electric field being normalized according to the relation
\begin{equation}
\overline{E}=E/E^{\dag }
\end{equation}

\subsection{Radial convection and diffusion}

From the conservative form of the Fokker-Planck equation in configuration
space
\begin{equation}
\frac{\partial f_{0}}{\partial t}\propto \frac{\partial }{\partial \psi }%
\left( \left| \nabla \psi \right| S_{\psi }\right) 
\end{equation}
where

\begin{equation}
S_{\psi }\propto -\Bbb{D}_{\psi }\left| \nabla \psi \right| \frac{\partial %
f_{0}}{\partial \psi }+\mathbf{F}_{\psi }f_{0}
\end{equation}

Introducing normalized coordinates, it turns out that
\begin{eqnarray}
\frac{\partial \overline{f}_{0}}{\partial \overline{t}} &\propto &\frac{1}{%
\nu _{e}^{\dag }\psi _{a}}\frac{\partial }{\partial \overline{\psi }}\left(
\left| \nabla \psi \right| S_{\psi }\right)   \nonumber \\
&=&\frac{1}{\nu _{e}^{\dag }\psi _{a}}\frac{\partial }{\partial \overline{%
\psi }}\left( \frac{\Bbb{D}_{\psi }}{\psi _{a}}\left| \nabla \psi \right|
^{2}\frac{\partial \overline{f}_{0}}{\partial \overline{\psi }}+\left|
\nabla \psi \right| \mathbf{F}_{\psi }\overline{f}_{0}\right)   \nonumber \\
&=&\frac{\partial }{\partial \overline{\psi }}\left( \frac{\Bbb{D}_{\psi }}{%
\nu _{e}^{\dag }\psi _{a}^{2}}\left| \nabla \psi \right| ^{2}\frac{\partial 
\overline{f}_{0}}{\partial \overline{\psi }}+\left| \nabla \psi \right| 
\frac{\mathbf{F}_{\psi }}{\nu _{e}^{\dag }\psi _{a}}\overline{f}_{0}\right) 
\end{eqnarray}
with
\begin{eqnarray}
\overline{\Bbb{D}}_{\psi } &=&\Bbb{D}_{\psi }/\Bbb{D}_{\psi }^{\dag } \\
\overline{\mathbf{F}}_{\psi } &=&\mathbf{F}_{\psi }/\mathbf{F}_{\psi }^{\dag
}
\end{eqnarray}
with
\begin{eqnarray}
\Bbb{D}_{\psi }^{\dag } &=&\nu _{e}^{\dag }\psi _{a}^{2} \\
\mathbf{F}_{\psi }^{\dag } &=&\nu _{e}^{\dag }\psi _{a}
\end{eqnarray}

However, since it is more convenient to handle diffusion and convection
coefficients that scale in $m^{2}/s$ and $m/s$ respectively,Radial 

\smallskip ***

\subsection{Cross-diffusion}

***

\subsection{Fluxes}

In momentum space, 
\begin{eqnarray}
S_{\mathbf{p}} &\propto &-\Bbb{D}_{\mathbf{p}}\frac{\partial f_{0}}{\partial %
p}+\mathbf{F}_{\mathbf{p}}f_{0}  \nonumber \\
&=&-\overline{\Bbb{D}}_{\mathbf{p}}\frac{\Bbb{D}_{\mathbf{p}}^{\dag
}n_{e}^{\dag }}{p_{th}^{\dag }}\frac{\partial \overline{f}_{0}}{\partial 
\overline{p}}+\overline{\mathbf{F}}_{\mathbf{p}}n_{e}^{\dag }\mathbf{F}_{%
\mathbf{p}}^{\dag }\overline{f}_{0}  \nonumber \\
&=&-\overline{\Bbb{D}}_{\mathbf{p}}\nu _{e}^{\dag }p_{th}^{\dag }n_{e}^{\dag
}\frac{\partial \overline{f}_{0}}{\partial \overline{p}}+\overline{\mathbf{F}%
}_{\mathbf{p}}n_{e}^{\dag }\nu _{e}^{\dag }p_{th}^{\dag }\overline{f}_{0} 
\nonumber \\
&=&\nu _{e}^{\dag }p_{th}^{\dag }n_{e}^{\dag }\left( -\overline{\Bbb{D}}_{%
\mathbf{p}}\frac{\partial \overline{f}_{0}}{\partial \overline{p}}+\overline{%
\mathbf{F}}_{\mathbf{p}}\overline{f}_{0}\right) 
\end{eqnarray}
and 
\begin{equation}
S_{\mathbf{p}}=\overline{S}_{\mathbf{p}}S_{\mathbf{p}}^{\dag }
\end{equation}
where
\begin{equation}
S_{\mathbf{p}}^{\dag }=\nu _{e}^{\dag }p_{th}^{\dag }n_{e}^{\dag
}=m_{e}c\beta _{th}^{\dag }\nu _{e}^{\dag }n_{e}^{\dag }
\end{equation}
and with these definitions, 
\[
\overline{S}_{\mathbf{p}}\propto -\overline{\Bbb{D}}_{\mathbf{p}}\frac{%
\partial \overline{f}_{0}}{\partial \overline{p}}+\overline{\mathbf{F}}_{%
\mathbf{p}}\overline{f}_{0}
\]

A similar procedure is applied for dynamics in configuration space, and 
\begin{equation}
S_{\psi }=\overline{S}_{\psi }S_{\psi }^{\dag }
\end{equation}
with
\begin{equation}
S_{\psi }^{\dag }=\nu _{e}^{\dag }\psi _{a}n_{e}^{\dag }
\end{equation}
and we get
\[
\]

\subsection{Stream function}

***

\subsection{Current density}

The current density is given by the relation
\begin{eqnarray}
J &\propto &2\pi q_{e}\int p^{2}dp\int_{-1}^{+1}\xi vf_{0}d\xi   \nonumber \\
&=&2\pi q_{e}\int \frac{p^{3}}{\gamma m_{e}}dp\int_{-1}^{+1}\xi f_{0}d\xi  
\nonumber \\
&=&2\pi \frac{q_{e}}{m_{e}}p_{th}^{\dag }n_{e}^{\dag }\int \frac{\overline{p}%
^{3}}{\gamma }d\overline{p}\int_{-1}^{+1}\xi \overline{f}_{0}d\xi 
\end{eqnarray}
since $f_{0}$ is normalized to $\overline{n}_{e},$ i.e.
\begin{equation}
2\pi \int \overline{p}^{2}d\overline{p}\int_{-1}^{+1}\overline{f}_{0}d\xi =%
\overline{n}_{e}
\end{equation}
Therefore,
\begin{equation}
\overline{J}=J/J^{\dag }
\end{equation}
where
\begin{equation}
J^{\dag }=\frac{q_{e}}{m_{e}}p_{th}^{\dag }n_{e}^{\dag }=q_{e}c\beta
_{th}^{\dag }n_{e}^{\dag }
\end{equation}

With this definition
\begin{equation}
\overline{J}\propto 2\pi \int \frac{\overline{p}^{3}}{\gamma }d\overline{p}%
\int_{-1}^{+1}\xi \overline{f}_{0}d\xi 
\end{equation}

\subsection{Power density}

The power density is deduced from the integral
\begin{eqnarray}
P &\propto &2\pi \int_{-1}^{+1}d\xi \int_{0}^{\infty }\frac{p^{3}}{\gamma
m_{e}}S_{\mathbf{p}}dp  \nonumber \\
&=&2\pi \frac{\nu _{e}^{\dag }p_{th}^{\dag 2}n_{e}^{\dag }}{m_{e}}%
\int_{-1}^{+1}d\xi \int_{0}^{\infty }\frac{\overline{p}^{3}}{\gamma }%
\overline{S}_{\mathbf{p}}d\overline{p}
\end{eqnarray}
and defining
\begin{equation}
\overline{P}=P/P^{\dag }
\end{equation}
one has
\begin{equation}
P^{\dag }=\frac{\nu _{e}^{\dag }p_{th}^{\dag 2}n_{e}^{\dag }}{m_{e}}%
=m_{e}c^{2}\beta _{th}^{\dag 2}\nu _{e}^{\dag }n_{e}^{\dag }
\end{equation}

With this definition
\begin{equation}
\overline{P}\propto 2\pi \int_{-1}^{+1}d\xi \int_{0}^{\infty }\frac{%
\overline{p}^{3}}{\gamma }\overline{S}_{\mathbf{p}}d\overline{p}
\end{equation}

\subsection{Eelctron runaway rate}

The electron runaway rate $\Gamma _{R}$ is given by the relation
\begin{eqnarray}
\Gamma _{R} &\propto &2\pi p^{2}\int_{-1}^{+1}S_{\mathbf{p}}d\xi   \nonumber
\\
&=&2\pi \overline{p}^{2}m_{e}c\beta _{th}^{\dag }\nu _{e}^{\dag }n_{e}^{\dag
}\int_{-1}^{+1}\overline{S}_{\mathbf{p}}d\xi 
\end{eqnarray}

Therefore, normalized expression is 
\begin{equation}
\overline{\Gamma }_{R}=\Gamma _{R}/\Gamma _{R}^{\dag }
\end{equation}
with
\begin{equation}
\Gamma _{R}^{\dag }=m_{e}c\beta _{th}^{\dag }\nu _{e}^{\dag }n_{e}^{\dag }
\end{equation}
so that 
\begin{equation}
\overline{\Gamma }_{R}\propto 2\pi \overline{p}^{2}\int_{-1}^{+1}\overline{S}%
_{\mathbf{p}}d\xi 
\end{equation}

The consistency of the normalized units may be benchmarked by the
approximate relation between the total current and the bulk current, the
difference coming from the electron runaway tail. For small $\Gamma _{R},$
the run-away distribution is independent of $v_{\Vert },$ so that $%
f_{0}\left( v_{\Vert }\gg v_{th}\right) \approx \left( \gamma /E\right)
\delta \left( v_{\bot }\right) .$ By integrating the current intergral out
to $v_{\max },$ one obtains
\begin{equation}
J\left( v_{\max }\right) \approx J_{bulk}+\frac{1}{2}\left( \frac{\Gamma _{R}%
}{E}\right) v_{\max }^{2}
\end{equation}
and in normalized units, 
\begin{equation}
\overline{J}\left( v_{\max }\right) J^{\dag }\approx \overline{J}%
_{bulk}J^{\dag }+\frac{1}{2}\left( \frac{\overline{\Gamma }_{R}\Gamma
_{R}^{\dag }}{\overline{E}E^{\dag }}\right) v_{th}^{\dag 2}\overline{v}%
_{\max }^{2}
\end{equation}
so that 
\begin{equation}
\overline{J}\left( v_{\max }\right) \approx \overline{J}_{bulk}+\frac{1}{2}%
\frac{v_{th}^{\dag 2}\Gamma _{R}^{\dag }}{E^{\dag }J^{\dag }}\left( \frac{%
\overline{\Gamma }_{R}}{\overline{E}}\right) \overline{v}_{\max }^{2}
\end{equation}
According to their respective expression, ***** A v\'{e}rifier **** 
\begin{eqnarray*}
\frac{v_{th}^{\dag 2}}{E^{\dag }J_{th}^{\dag }} &=&\frac{c^{2}\beta
_{th}^{\dag 2}m_{e}c\beta _{th}^{\dag }\nu _{e}^{\dag }n_{e}^{\dag }}{
q_{e}c\beta _{th}^{\dag }n_{e}^{\dag }m_{e}c\beta _{th}^{\dag }\nu
_{e}^{\dag }/q_{e}} \\
&=&
\end{eqnarray*}

\subsection{Bremsstrahlung emission}

\subsection{Units}

In the code, the following units are used

\[
\begin{tabular}{ll}
Quantity & Units \\ 
Temperature & $keV$ \\ 
Density & $m^{-3}$ \\ 
Time & $s$ \\ 
Frequency & $s^{-1}$%
\end{tabular}
\]
